\documentclass[11pt]{article}
\usepackage{fontspec,amsmath,amssymb,amsthm,hyperref}
\usepackage[margin=1in]{geometry}
\setlength{\emergencystretch}{3em}
\title{Construct the physical particle drift in a coordinate L\ensuremath{{}^{2}} space}
\author{LangevinLean mathematical exposition}
\date{}
\begin{document}\maketitle
\noindent Source module: \texttt{ProximalLangevin/MeanFieldParticleDrift.lean}. Compile this proof document with XeLaTeX or LuaLaTeX. The accompanying HTML places the complete Lean declarations beside every proof.


Let \(H_N=L^2(P;\Theta)\times\prod_{i<N}L^2(P;X)\) with norm \(\|z\|=\max(\|z_\theta\|_2,\max_i\|z_i\|_2)\). This Banach norm controls every tagged coordinate directly, with no exchangeability hypothesis.


\section{Coordinate inclusions and nonexpansive projections}


Define the coordinate product state and the latent inclusion \(I_x(x)=(0,x)\). The Euclidean square proves \(\|I_xx\|=\|x\|\), while the latent projection has norm at most one. For any bounded linear pointwise contraction \(C\), integrating \(\|C(X)\|^2\le\|X\|^2\) proves \(\|C\circ X\|_2\le\|X\|_2\).



\noindent Lean declarations: \texttt{MeanFieldParticleState}, \texttt{jointLatentInclusionCLM}, \texttt{jointLatentInclusionCLM\_norm}, \texttt{norm\_jointLatentCLM\_apply\_le}, \texttt{lp\_linear\_norm\_le}.

\section{Assemble one actual joint random pair}


Set \(J(\Theta,X)=I_\theta\Theta+I_xX\). Its representative is the genuine pair \((\Theta(\omega),X(\omega))\) almost everywhere. Linearity of the inclusions and their norm bounds give

\[\|J(\Theta,X)-J(\Theta',X')\|_2\le\|\Theta-\Theta'\|_2+\|X-X'\|_2.\]

\noindent Lean declarations: \texttt{lpJointPair}, \texttt{lpJointPair\_coe}, \texttt{lpJointPair\_difference\_le}.

\section{Define the gradient and empirical particle drift}


Apply the actual gradient \(\nabla V\) to each pair \(J(\Theta,X_i)\). The parameter force is the empirical average and each latent force is its own gradient block:

\[\mathcal A_N(z)=\left(\frac1N\sum_i\nabla_\theta V(\Theta,X_i),\ (\nabla_xV(\Theta,X_i))_i\right).\]

Almost everywhere composition and finite-sum identities verify these exact coordinate formulas. The definition does not replace the empirical parameter force by a mean-field expectation.



\noindent Lean declarations: \texttt{lpParticleGradient}, \texttt{lpParticleDrift}, \texttt{lpParticleGradient\_coe}, \texttt{lpParticleDrift\_parameter\_coe}, \texttt{lpParticleDrift\_latent\_coe}.

\section{Dimension-free Lipschitz control constructs the actual solution}


Each coordinate difference is bounded by the product-state norm. The joint-pair estimate therefore gives \(\|\nabla V(J_i z)-\nabla V(J_i w)\|_2\le2L\|z-w\|\). Averaging \(N\) parameter terms and dividing by \(N>0\) preserves this bound; the maximum of the latent terms also preserves it. Hence \(\operatorname{Lip}(\mathcal A_N)\le2L\), independently of \(N\). Picard constructs continuous \(Z\) satisfying \(Z_t-Z_0=-\int_0^t\mathcal A_N(Z_s)ds+B_t\) for any continuous input with \(B_0=0\).



\noindent Lean declarations: \texttt{lpParticleGradient\_difference\_le}, \texttt{lpParticleDrift\_lipschitz}, \texttt{exists\_lp\_particle\_integral}.

\end{document}
